\chapter{Chapter06 — Small Bias Variants}
%%% generated by the blueprint pipeline from FABL.Chapter06.Constructions.SmallBiasVariants
\section{Overview}

Declarations of the Lean module \texttt{FABL.Chapter06.Constructions.SmallBiasVariants}.

\section{Declarations}

\begin{definition}[Field pairing coefficient]
\lean{FABL.fieldPairingCoefficient}
\label{FABL.fieldPairingCoefficient}
\leanok
\uses{FABL.BinaryExtensionField, FABL.F₂Cube}
Let $$ be a binary vector of length $d$, and let $a$ assign an element of the binary
extension field of degree $$ to each index in $$. Define their field pairing
coefficient to be the sum over all indices $i$ of the scalar multiple of $a_i$ by the
binary coefficient $_i$.
\end{definition}

\begin{definition}[Field pairing generator]
\lean{FABL.fieldPairingGenerator}
\statementsource{Analysis-of-Boolean-Functions-by-Ryan-ODonnell}{pdf-ffbe9c7b3306-p153-b003}
\label{FABL.fieldPairingGenerator}
\leanok
\uses{FABL.BinaryExtensionField, FABL.F₂Cube, FABL.binaryExtensionEncode, FABL.f₂DotProduct}
Let $d,\ell$ be natural numbers with $\ell\ne0$, let $a_0,\ldots,a_{d-1}$ be elements of
the binary extension field of degree $\ell$, and let $s$ be another element of that
field. Define a binary vector of length $d$ whose $i$-th coordinate is the binary dot
product of the binary encodings of $a_i$ and $s$.
\end{definition}

\begin{theorem}[Pairing through a binary extension field]
\lean{FABL.f₂DotProduct_fieldPairingGenerator}
\leanok
\uses{FABL.BinaryExtensionField, FABL.F₂Cube, FABL.binaryExtensionEncode, FABL.fieldPairingCoefficient, FABL.fieldPairingGenerator, FABL.f₂DotProduct}
\difficulty{4}
Let $\ell$ be a nonzero integer, and let $d$ be a nonnegative integer. Use a basis of
the degree-$\ell$ binary extension field to encode its elements as binary vectors of
length $\ell$. For $\gamma\in\bbr_2^d$, a family $a_0,\ldots,a_{d-1}$ in that field, and
an element $s$ of the field, the dot product of $\gamma$ with the vector whose $i$th
coordinate pairs the encodings of $a_i$ and $s$ equals the dot product of the encodings
of $\sum_{i=0}^{d-1}\gamma_i a_i$ and $s$.
\end{theorem}

\begin{theorem}[Walsh characters and field pairings]
\lean{FABL.vectorWalshCharacter_fieldPairingGenerator}
\label{FABL.vectorWalshCharacter_fieldPairingGenerator}
\leanok
\uses{FABL.BinaryExtensionField, FABL.F₂Cube, FABL.binaryExtensionEncode, FABL.fieldPairingCoefficient, FABL.fieldPairingGenerator, FABL.f₂DotProduct_fieldPairingGenerator, FABL.vectorWalshCharacter, FABL.vectorWalshCharacter_apply}
\difficulty{3}
Let $d,n,\ell,t$ be natural numbers with $\ell\ne0$. Let $\gamma\in\mathbb F_2^d$, let
$a=(a_i)_{i=0}^{d-1}$ be elements of the binary extension field of degree $\ell$, and
let $s$ be an element of that field. Choose a basis to identify the field with $\mathbb
F_2^\ell$, and let the field pairing generator have $i$th coordinate the binary dot
product of the coordinate vectors of $a_i$ and $s$. The Walsh character indexed by
$\gamma$ evaluated at this generated vector equals the Walsh character indexed by the
coordinate vector of $\sum_{i=0}^{d-1}\gamma_i a_i$, evaluated at the coordinate vector
of $s$.
\end{theorem}

\begin{theorem}[Walsh average of a field pairing generator]
\lean{FABL.expect_fieldPairingGenerator}
\proofsource{Analysis-of-Boolean-Functions-by-Ryan-ODonnell}{pdf-ffbe9c7b3306-p153-b004}
\label{FABL.expect_fieldPairingGenerator}
\leanok
\uses{FABL.BinaryExtensionField, FABL.F₂Cube, FABL.binaryExtensionEncode, FABL.expect_vectorWalshCharacter, FABL.fieldPairingCoefficient, FABL.fieldPairingGenerator, FABL.vectorWalshCharacter, FABL.vectorWalshCharacter_fieldPairingGenerator}
\difficulty{4}
Let $\ell\neq 0$, let $\mathbb F$ be the degree-$\ell$ extension field of $\mathbb F_2$,
let $\gamma\in\mathbb F_2^d$, and let $a_1,\ldots,a_d\in\mathbb F$. Choose $s$ uniformly
from $\mathbb F$ and form the sign vector whose $i$th coordinate is
$(-1)^{\operatorname{Tr}_{\mathbb F/\mathbb F_2}(a_i s)}$. The average of its Walsh
character indexed by $\gamma$ is $1$ if $\sum_{i=1}^d\gamma_i a_i=0$ in $\mathbb F$, and
$0$ otherwise.
\end{theorem}

\begin{theorem}[Character expectation and polynomial roots]
\lean{FABL.fieldPairingGenerator_characterExpectation_eq_rootProbability}
\proofsource{Analysis-of-Boolean-Functions-by-Ryan-ODonnell}{pdf-ffbe9c7b3306-p153-b004}
\label{FABL.fieldPairingGenerator_characterExpectation_eq_rootProbability}
\leanok
\uses{FABL.BinaryExtensionField, FABL.F₂Cube, FABL.binaryExtensionField_natCard, FABL.expect_fieldPairingGenerator, FABL.fieldPairingCoefficient, FABL.fieldPairingGenerator, FABL.vectorWalshCharacter}
\difficulty{5}
Let $\ell>0$, let $\gamma$ be a $d$-coordinate binary vector, let $a$ map each element
of $\mathbb F_{2^\ell}$ to a $d$-tuple over that field, and let $p$ be a nonzero
polynomial over $\mathbb F_{2^\ell}$. Suppose $p(r)=\sum_i\gamma_i a(r)_i$ for every
field element $r$. The average of the Walsh character indexed by $\gamma$ on the
generated vector, when its two field inputs are chosen independently and uniformly,
equals the number of roots of $p$ in $\mathbb F_{2^\ell}$ divided by $2^\ell$.
\end{theorem}

\begin{definition}[Field pairing density]
\lean{FABL.fieldPairingDensity}
\statementsource{Analysis-of-Boolean-Functions-by-Ryan-ODonnell}{pdf-ffbe9c7b3306-p153-b003}
\label{FABL.fieldPairingDensity}
\leanok
\uses{FABL.BinaryExtensionField, FABL.ProbabilityDensity, FABL.ProbabilityDensity.uniformPushforward, FABL.fieldPairingGenerator}
Let $d,\ell$ be natural numbers with $\ell\ne0$, and for each element $r$ of the binary
extension field of degree $\ell$, let $a(r)$ be a length-$d$ list of elements of that
field. Define a probability density on the binary cube of dimension $d$ by the uniform
pushforward of the field pairing generator applied to a uniformly chosen pair $(r,s)$ of
field elements, using $a(r)$ and $s$.
\end{definition}

\begin{theorem}[Fourier coefficient and polynomial root probability]
\lean{FABL.vectorFourierCoeff_fieldPairingDensity_eq_rootProbability}
\proofsource{Analysis-of-Boolean-Functions-by-Ryan-ODonnell}{pdf-ffbe9c7b3306-p153-b004}
\label{FABL.vectorFourierCoeff_fieldPairingDensity_eq_rootProbability}
\leanok
\uses{FABL.BinaryExtensionField, FABL.F₂Cube, FABL.ProbabilityDensity.expectation, FABL.ProbabilityDensity.expectation_uniformPushforward, FABL.fieldPairingCoefficient, FABL.fieldPairingDensity, FABL.fieldPairingGenerator_characterExpectation_eq_rootProbability, FABL.vectorFourierCoeff, FABL.vectorFourierCoeff_eq_expect, FABL.vectorWalshCharacter}
\difficulty{3}
Let $\ell\ne0$, let $\gamma$ be a $d$-coordinate binary vector, and let $a$ assign to
each element $r$ of the degree-$\ell$ extension field of $\mathbb{F}_2$ a $d$-tuple of
field elements. Let $p$ be a nonzero polynomial over that field whose value at every $r$
equals the field sum of the coordinates of $\gamma$ multiplied by the corresponding
entries of $a(r)$. Then the vector Fourier coefficient of the density generated by the
field pairing is the number of roots of $p$ in the extension field divided by $2^\ell$.
\end{theorem}

\begin{theorem}[Bias from polynomial field pairings]
\lean{FABL.fieldPairingDensity_isBiased_of_polynomial}
\proofsource{Analysis-of-Boolean-Functions-by-Ryan-ODonnell}{pdf-ffbe9c7b3306-p153-b004}
\label{FABL.fieldPairingDensity_isBiased_of_polynomial}
\leanok
\uses{FABL.BinaryExtensionField, FABL.F₂Cube, FABL.ProbabilityDensity.IsBiased, FABL.fieldPairingCoefficient, FABL.fieldPairingDensity, FABL.ncard_rootSet_le_natDegree, FABL.vectorFourierCoeff_fieldPairingDensity_eq_rootProbability}
\difficulty{3}
Let $\ell>0$, let $a$ assign to each element $r$ of the binary extension field $\mathbb
F_{2^\ell}$ a vector of $d$ elements of that field, and let $p$ assign to each
$d$-coordinate binary vector a polynomial over $\mathbb F_{2^\ell}$. Suppose every
polynomial $p(\gamma)$ with $\gamma\ne0$ is nonzero, its value at every $r$ equals
$\sum_i\gamma_i a(r)_i$, and its degree is at most $D$. The probability density obtained
by choosing two field elements independently and uniformly and generating the binary
vector from their pairings with $a$ is biased: every nonzero Fourier coefficient has
absolute value at most $D/2^\ell$.
\end{theorem}

\begin{theorem}[Dyadic field ratio]
\lean{FABL.dyadicFieldRatio_eq_invPow}
\proofsource{Analysis-of-Boolean-Functions-by-Ryan-ODonnell}{pdf-ffbe9c7b3306-p153-b004}
\label{FABL.dyadicFieldRatio_eq_invPow}
\leanok
\difficulty{4}
For natural numbers $t$ and $\ell$ with $t\leq\ell$, the ratio $2^{\ell-t}/2^\ell$
equals $2^{-t}$.
\end{theorem}

\begin{definition}[Shifted field family]
\lean{FABL.shiftedSmallBiasFieldFamily}
\statementsource{Analysis-of-Boolean-Functions-by-Ryan-ODonnell}{pdf-ffbe9c7b3306-p166-b005}
\label{FABL.shiftedSmallBiasFieldFamily}
\leanok
\uses{FABL.BinaryExtensionField}
For natural numbers $n,\ell$, a field element $r$ in the binary extension field of
degree $\ell$ determines the length-$n$ family whose entry at index $i$ is $r^i$.
\end{definition}

\begin{definition}[Shifted small bias polynomial]
\lean{FABL.shiftedSmallBiasPolynomial}
\statementsource{Analysis-of-Boolean-Functions-by-Ryan-ODonnell}{pdf-ffbe9c7b3306-p166-b005}
\label{FABL.shiftedSmallBiasPolynomial}
\leanok
\uses{FABL.BinaryExtensionField, FABL.F₂Cube, FABL.𝔽₂}
For natural numbers $n,\ell$ and a vector $\gamma\in(\mathbb{Z}/2\mathbb{Z})^n$, define
a polynomial over the binary extension field of degree $\ell$ by
$\sum_{i=0}^{n-1}\gamma_iX^i$, where each coefficient is mapped into that field by its
canonical field embedding.
\end{definition}

\begin{theorem}[Evaluation of the shifted polynomial]
\lean{FABL.shiftedSmallBiasPolynomial_eval}
\label{FABL.shiftedSmallBiasPolynomial_eval}
\leanok
\uses{FABL.BinaryExtensionField, FABL.F₂Cube, FABL.fieldPairingCoefficient, FABL.shiftedSmallBiasFieldFamily, FABL.shiftedSmallBiasPolynomial}
\difficulty{1}
Let $\ell,n\in\mathbb{N}$, let $\mathbb{F}$ be the field with $2^\ell$ elements, let
$\gamma\in\{0,1\}^n$, and let $r\in\mathbb{F}$. The polynomial $\sum_{i=0}^{n-1}\gamma_i
X^i$, with each $\gamma_i$ viewed in $\mathbb{F}$, evaluates at $r$ to
$\sum_{i=0}^{n-1}\gamma_i r^i$.
\end{theorem}

\begin{theorem}[Nonvanishing of the shifted polynomial]
\lean{FABL.shiftedSmallBiasPolynomial_ne_zero}
\label{FABL.shiftedSmallBiasPolynomial_ne_zero}
\leanok
\uses{FABL.BinaryExtensionField, FABL.F₂Cube, FABL.shiftedSmallBiasPolynomial, FABL.𝔽₂}
\difficulty{4}
Let $\ell,n\in\mathbb{N}$ and let $\gamma\in\{0,1\}^n$ be nonzero. Then the polynomial
$\sum_{i=0}^{n-1}\gamma_i X^i$, with coefficients viewed in the field with $2^\ell$
elements, is nonzero.
\end{theorem}

\begin{theorem}[Shifted small bias polynomial nat degree le]
\lean{FABL.shiftedSmallBiasPolynomial_natDegree_le}
\statementsource{Analysis-of-Boolean-Functions-by-Ryan-ODonnell}{pdf-ffbe9c7b3306-p166-b005}
\label{FABL.shiftedSmallBiasPolynomial_natDegree_le}
\leanok
\uses{FABL.F₂Cube, FABL.shiftedSmallBiasPolynomial}
\difficulty{3}
% dataset-pipeline: stub -- prose pending build_informal_statements.py
Statement pending regeneration.
\end{theorem}

\begin{definition}[Shifted small bias generator]
\lean{FABL.shiftedSmallBiasGenerator}
\statementsource{Analysis-of-Boolean-Functions-by-Ryan-ODonnell}{pdf-ffbe9c7b3306-p166-b005}
\label{FABL.shiftedSmallBiasGenerator}
\leanok
\uses{FABL.BinaryExtensionField, FABL.F₂Cube, FABL.fieldPairingGenerator, FABL.shiftedSmallBiasFieldFamily}
Let $n,\ell$ be natural numbers with $\ell\ne0$, and let $r,s$ be elements of the binary
extension field of degree $\ell$. Define a binary vector of length $n$ by applying the
field pairing generator to $s$ and the family whose $i$-th entry is $r^i$.
\end{definition}

\begin{theorem}[Character expectation and polynomial roots]
\lean{FABL.shiftedSmallBiasGenerator_characterExpectation_eq_rootProbability}
\proofsource{Analysis-of-Boolean-Functions-by-Ryan-ODonnell}{pdf-ffbe9c7b3306-p153-b004}
\label{FABL.shiftedSmallBiasGenerator_characterExpectation_eq_rootProbability}
\leanok
\uses{FABL.BinaryExtensionField, FABL.F₂Cube, FABL.fieldPairingGenerator_characterExpectation_eq_rootProbability, FABL.shiftedSmallBiasFieldFamily, FABL.shiftedSmallBiasGenerator, FABL.shiftedSmallBiasPolynomial, FABL.shiftedSmallBiasPolynomial_eval, FABL.shiftedSmallBiasPolynomial_ne_zero, FABL.vectorWalshCharacter}
\difficulty{1}
Let $n,\ell$ be natural numbers with $\ell\ne0$, and let $\gamma\in\mathbb F_2^n$ be
nonzero. Choose $r,s$ independently and uniformly from the binary extension field
$\mathbb F_{2^\ell}$. The expectation of the Walsh character indexed by $\gamma$,
evaluated at the shifted small bias generator determined by $r,s$, equals the number of
roots in $\mathbb F_{2^\ell}$ of $p_\gamma(X)=\sum_{i=0}^{n-1}\gamma_iX^i$, divided by
$2^\ell$.
\end{theorem}

\begin{definition}[Shifted small bias density]
\lean{FABL.shiftedSmallBiasGeneratorDensity}
\statementsource{Analysis-of-Boolean-Functions-by-Ryan-ODonnell}{pdf-ffbe9c7b3306-p166-b005}
\label{FABL.shiftedSmallBiasGeneratorDensity}
\leanok
\uses{FABL.ProbabilityDensity, FABL.fieldPairingDensity, FABL.shiftedSmallBiasFieldFamily}
For a natural number $n$ and a nonzero natural number $\ell$, the shifted small bias
generator density is the probability density on $\mathbb{F}_2^n$ obtained by pairing the
shifted small bias field family with the field pairing construction.
\end{definition}

\begin{theorem}[Bias of the shifted small bias generator]
\lean{FABL.shiftedSmallBiasGeneratorDensity_isBiased}
\proofsource{Analysis-of-Boolean-Functions-by-Ryan-ODonnell}{
  pdf-ffbe9c7b3306-p166-b005,
  pdf-ffbe9c7b3306-p153-b004
}
\label{FABL.shiftedSmallBiasGeneratorDensity_isBiased}
\leanok
\uses{FABL.ProbabilityDensity.IsBiased, FABL.fieldPairingDensity_isBiased_of_polynomial, FABL.shiftedSmallBiasFieldFamily, FABL.shiftedSmallBiasGeneratorDensity, FABL.shiftedSmallBiasPolynomial, FABL.shiftedSmallBiasPolynomial_eval, FABL.shiftedSmallBiasPolynomial_natDegree_le, FABL.shiftedSmallBiasPolynomial_ne_zero}
\difficulty{1}
Let $n,\ell$ be natural numbers with $\ell\ne0$. The probability density on $\mathbb
F_2^n$ obtained by choosing the shifted small bias generator uniformly from its
parameter space is biased with parameter $(n-1)/2^\ell$: for every nonzero
$\gamma\in\mathbb F_2^n$, the absolute value of its vector Fourier coefficient is at
most $(n-1)/2^\ell$.
\end{theorem}

\begin{theorem}[Dyadic parameter identity]
\lean{FABL.shiftedSmallBias_dyadicParameter}
\label{FABL.shiftedSmallBias_dyadicParameter}
\leanok
\uses{FABL.dyadicFieldRatio_eq_invPow}
\difficulty{4}
Let $t,\ell,n\in\mathbb{N}$ satisfy $t\leq\ell$ and $n=2^{\ell-t}$. Then, as an identity
in $\mathbb{R}$, $\frac{n-1}{2^\ell}=\left(\frac12\right)^t-\left(\frac12\right)^\ell$.
\end{theorem}

\begin{theorem}[Dyadic bias of the shifted generator]
\lean{FABL.shiftedSmallBiasGeneratorDensity_isBiased_dyadic}
\proofsource{Analysis-of-Boolean-Functions-by-Ryan-ODonnell}{pdf-ffbe9c7b3306-p166-b005}
\label{FABL.shiftedSmallBiasGeneratorDensity_isBiased_dyadic}
\leanok
\uses{FABL.ProbabilityDensity.IsBiased, FABL.shiftedSmallBiasGeneratorDensity, FABL.shiftedSmallBiasGeneratorDensity_isBiased, FABL.shiftedSmallBias_dyadicParameter}
\difficulty{1}
Let $n,\ell,t$ be natural numbers such that $\ell\ne0$, $t\le\ell$, and $n=2^{\ell-t}$.
The probability density on $\mathbb F_2^n$ obtained by choosing the shifted small bias
generator uniformly from its parameter space is biased with parameter
$2^{-t}-2^{-\ell}$: every nonzero vector Fourier coefficient has absolute value at most
this quantity.
\end{theorem}

\begin{definition}[Binary extension basis vector]
\lean{FABL.binaryExtensionBasisVector}
\label{FABL.binaryExtensionBasisVector}
\leanok
\uses{FABL.BinaryExtensionField, FABL.binaryExtensionEncode}
Let $$ be a nonzero natural number and let $j$ be an index from $0$ to $-1$. Define
the corresponding basis vector in the binary extension field of degree $$ as the
inverse image, under the binary extension encoding, of the coordinate vector that is $1$
at $j$ and $0$ elsewhere.
\end{definition}

\begin{theorem}[Encoding a basis vector]
\lean{FABL.binaryExtensionEncode_basisVector}
\label{FABL.binaryExtensionEncode_basisVector}
\leanok
\uses{FABL.binaryExtensionBasisVector, FABL.binaryExtensionEncode}
\difficulty{1}
Let $\ell$ be a positive integer, and let $j$ be a coordinate in $\{0,\ldots,\ell-1\}$.
Under the linear encoding of the degree-$\ell$ binary extension field by its coordinates
in a basis, the field element corresponding to the $j$th basis vector is sent to the
coordinate vector that is $1$ at $j$ and $0$ elsewhere.
\end{theorem}

\begin{definition}[Grouped basis coefficient]
\lean{FABL.basisExpandedGroupedCoefficient}
\label{FABL.basisExpandedGroupedCoefficient}
\leanok
\uses{FABL.BinaryExtensionField, FABL.F₂Cube, FABL.binaryExtensionEncode}
Let $$ be nonzero, let $$ be a binary vector of length $n$, and let $i$ be an index
from $0$ to $n-1$. Define the grouped coefficient in the binary extension field of
degree $$ to be the inverse image, under the binary extension encoding, of the
length-$$ vector whose $j$th coordinate is the coordinate of $$ indexed by the pair
$(i,j)$.
\end{definition}

\begin{theorem}[Grouped coefficient basis expansion]
\lean{FABL.basisExpandedGroupedCoefficient_eq_sum}
\label{FABL.basisExpandedGroupedCoefficient_eq_sum}
\leanok
\uses{FABL.F₂Cube, FABL.basisExpandedGroupedCoefficient, FABL.binaryExtensionBasisVector, FABL.binaryExtensionEncode, FABL.binaryExtensionEncode_basisVector}
\difficulty{3}
Let $\ell$ be a nonzero natural number, let $\gamma$ assign an element of $\mathbb
F_2$ to each pair of indices in $\mathrm{Fin}(n)\times\mathrm{Fin}(\ell)$, and
choose $i\in\mathrm{Fin}(n)$. The element of the degree-$\ell$ binary extension field
encoded by the coefficients $\gamma(i,j)$ is the sum, over $j\in\mathrm{Fin}(\ell)$,
of $\gamma(i,j)$ times the $j$th vector of the chosen coordinate basis.
\end{theorem}

\begin{theorem}[A nonzero grouped extension coefficient]
\lean{FABL.exists_basisExpandedGroupedCoefficient_ne_zero}
\label{FABL.exists_basisExpandedGroupedCoefficient_ne_zero}
\leanok
\uses{FABL.F₂Cube, FABL.basisExpandedGroupedCoefficient, FABL.binaryExtensionEncode}
\difficulty{5}
Let $\ell>0$, and let $\gamma$ be a nonzero binary string of length $n\ell$, viewed as
$n$ consecutive blocks of length $\ell$. At least one block, interpreted through a fixed
basis of the degree-$\ell$ binary extension field, represents a nonzero field element.
\end{theorem}

\begin{definition}[Basis expanded field family]
\lean{FABL.basisExpandedSmallBiasFieldFamily}
\statementsource{Analysis-of-Boolean-Functions-by-Ryan-ODonnell}{pdf-ffbe9c7b3306-p166-b005}
\label{FABL.basisExpandedSmallBiasFieldFamily}
\leanok
\uses{FABL.BinaryExtensionField, FABL.binaryExtensionBasisVector}
For natural numbers $n$ and nonzero $\ell$, and an element $r$ of the binary extension
field of degree $\ell$, define a family indexed by $q\in\mathrm{Fin}(n\ell)$ as follows.
Write the index corresponding to $q$ as
$(i,j)\in\mathrm{Fin}(n)\times\mathrm{Fin}(\ell)$; its family value is the $j$th binary
extension basis vector multiplied by $r^{i+1}$.
\end{definition}

\begin{definition}[Basis expanded polynomial]
\lean{FABL.basisExpandedSmallBiasPolynomial}
\statementsource{Analysis-of-Boolean-Functions-by-Ryan-ODonnell}{pdf-ffbe9c7b3306-p166-b005}
\label{FABL.basisExpandedSmallBiasPolynomial}
\leanok
\uses{FABL.BinaryExtensionField, FABL.F₂Cube, FABL.basisExpandedGroupedCoefficient}
For a nonzero natural number $\ell$ and a binary cube $\gamma\in\mathbb{F}_2^{n\ell}$,
define a polynomial over the binary extension field of degree $\ell$ by
$\sum_{i=0}^{n-1} c_i X^{i+1}$, where $c_i$ is the grouped coefficient associated with
$\gamma$ and index $i$.
\end{definition}

\begin{theorem}[Basis expanded small bias polynomial eval]
\lean{FABL.basisExpandedSmallBiasPolynomial_eval}
\statementsource{Analysis-of-Boolean-Functions-by-Ryan-ODonnell}{pdf-ffbe9c7b3306-p166-b005}
\label{FABL.basisExpandedSmallBiasPolynomial_eval}
\leanok
\uses{FABL.BinaryExtensionField, FABL.F₂Cube, FABL.basisExpandedGroupedCoefficient, FABL.basisExpandedGroupedCoefficient_eq_sum, FABL.basisExpandedSmallBiasFieldFamily, FABL.basisExpandedSmallBiasPolynomial, FABL.binaryExtensionBasisVector, FABL.fieldPairingCoefficient}
\difficulty{5}
% dataset-pipeline: stub -- prose pending build_informal_statements.py
Statement pending regeneration.
\end{theorem}

\begin{theorem}[Nonvanishing of the expanded polynomial]
\lean{FABL.basisExpandedSmallBiasPolynomial_ne_zero}
\label{FABL.basisExpandedSmallBiasPolynomial_ne_zero}
\leanok
\uses{FABL.F₂Cube, FABL.basisExpandedGroupedCoefficient, FABL.basisExpandedSmallBiasPolynomial, FABL.exists_basisExpandedGroupedCoefficient_ne_zero}
\difficulty{4}
Let $n$ and $\ell$ be natural numbers with $\ell \ne 0$, and let $\gamma$ be a nonzero
vector in $(\mathbb{F}_2)^{n\ell}$. Group its coordinates into $n$ blocks of length
$\ell$, and use the chosen basis of the degree-$\ell$ binary extension field to
interpret each block as a coefficient $a_i$ in that field. Then the polynomial
$\sum_{i=0}^{n-1} a_i X^{i+1}$ over the binary extension field is nonzero.
\end{theorem}

\begin{theorem}[Degree bound for the encoded polynomial]
\lean{FABL.basisExpandedSmallBiasPolynomial_natDegree_le}
\label{FABL.basisExpandedSmallBiasPolynomial_natDegree_le}
\leanok
\uses{FABL.F₂Cube, FABL.basisExpandedSmallBiasPolynomial}
\difficulty{3}
Let $n$ and $ell$ be natural numbers with $ell \ne 0$, and let $\gamma$ be a function
from the $nell$ coordinates to $\mathbb{F}_2$. Form a polynomial over the degree-$ell$
extension field of $\mathbb{F}_2$ by grouping the coordinates of $\gamma$ into $n$
blocks of length $ell$, interpreting each block in the extension field using its chosen
basis, and using these values as the coefficients of $X, X^2, \ldots, X^n$,
respectively. The natural degree of this polynomial is at most $n$.
\end{theorem}

\begin{definition}[Basis expanded generator]
\lean{FABL.basisExpandedSmallBiasGenerator}
\statementsource{Analysis-of-Boolean-Functions-by-Ryan-ODonnell}{pdf-ffbe9c7b3306-p166-b005}
\label{FABL.basisExpandedSmallBiasGenerator}
\leanok
\uses{FABL.BinaryExtensionField, FABL.F₂Cube, FABL.basisExpandedSmallBiasFieldFamily, FABL.fieldPairingGenerator}
For natural numbers $n$ and nonzero $\ell$, and elements $r,s$ of the binary extension
field of degree $\ell$, the basis expanded small bias generator is the binary cube of
length $n\ell$ obtained by applying the field pairing generator to the basis expanded
small bias field family at $r$, with second parameter $s$.
\end{definition}

\begin{theorem}[Generator coordinate formula]
\lean{FABL.basisExpandedSmallBiasGenerator_apply_finProdFinEquiv}
\proofsource{Analysis-of-Boolean-Functions-by-Ryan-ODonnell}{pdf-ffbe9c7b3306-p166-b005}
\label{FABL.basisExpandedSmallBiasGenerator_apply_finProdFinEquiv}
\leanok
\uses{FABL.BinaryExtensionField, FABL.basisExpandedSmallBiasFieldFamily, FABL.basisExpandedSmallBiasGenerator, FABL.binaryExtensionBasisVector, FABL.binaryExtensionEncode, FABL.fieldPairingGenerator, FABL.f₂DotProduct}
\difficulty{1}
Let $n$ be a natural number and let $[?]$ be a positive integer. In the binary
extension field $K$ of degree $[?]$ over $[?]$, let $E:K\to\mathbb{F}_2^{\ell}$ be the
coordinate isomorphism and let $b_j$ be the basis vector corresponding to $j$. For every
$r,s\in K$, $i\in\{0,\ldots,n-1\}$, and $j\in\{0,\ldots,\ell-1\}$, the $(i,j)$
coordinate of the basis-expanded generator is the $[?]$ dot product of $E(b_j r^{i+1})$
and $E(s)$.
\end{theorem}

\begin{theorem}[Basis expanded small bias generator character expectation eq root probability]
\lean{FABL.basisExpandedSmallBiasGenerator_characterExpectation_eq_rootProbability}
\statementsource{Analysis-of-Boolean-Functions-by-Ryan-ODonnell}{
  pdf-ffbe9c7b3306-p153-b004,
  pdf-ffbe9c7b3306-p166-b005
}
\label{FABL.basisExpandedSmallBiasGenerator_characterExpectation_eq_rootProbability}
\leanok
\uses{FABL.BinaryExtensionField, FABL.F₂Cube, FABL.basisExpandedSmallBiasFieldFamily, FABL.basisExpandedSmallBiasGenerator, FABL.basisExpandedSmallBiasPolynomial, FABL.basisExpandedSmallBiasPolynomial_eval, FABL.basisExpandedSmallBiasPolynomial_ne_zero, FABL.fieldPairingGenerator_characterExpectation_eq_rootProbability, FABL.vectorWalshCharacter}
\difficulty{1}
% dataset-pipeline: stub -- prose pending build_informal_statements.py
Statement pending regeneration.
\end{theorem}

\begin{definition}[Basis expanded generator density]
\lean{FABL.basisExpandedSmallBiasGeneratorDensity}
\statementsource{Analysis-of-Boolean-Functions-by-Ryan-ODonnell}{pdf-ffbe9c7b3306-p166-b005}
\label{FABL.basisExpandedSmallBiasGeneratorDensity}
\leanok
\uses{FABL.ProbabilityDensity, FABL.basisExpandedSmallBiasFieldFamily, FABL.fieldPairingDensity}
For natural numbers $n$ and nonzero $\ell$, the basis expanded small bias generator
density is the probability density on $\mathbb{F}_2^{n\ell}$ obtained by applying the
field pairing density construction to the basis expanded small bias field family.
\end{definition}

\begin{theorem}[Bias of the generator density]
\lean{FABL.basisExpandedSmallBiasGeneratorDensity_isBiased}
\proofsource{Analysis-of-Boolean-Functions-by-Ryan-ODonnell}{
  pdf-ffbe9c7b3306-p166-b005,
  pdf-ffbe9c7b3306-p153-b004
}
\label{FABL.basisExpandedSmallBiasGeneratorDensity_isBiased}
\leanok
\uses{FABL.ProbabilityDensity.IsBiased, FABL.basisExpandedSmallBiasFieldFamily, FABL.basisExpandedSmallBiasGeneratorDensity, FABL.basisExpandedSmallBiasPolynomial, FABL.basisExpandedSmallBiasPolynomial_eval, FABL.basisExpandedSmallBiasPolynomial_natDegree_le, FABL.basisExpandedSmallBiasPolynomial_ne_zero, FABL.fieldPairingDensity_isBiased_of_polynomial}
\difficulty{1}
Let $n$ and $\ell$ be natural numbers with $\ell\ne 0$. The probability density
on $\mathbb{F}_2^{n\ell}$ obtained by choosing two elements of the binary extension
field of degree $\ell$ independently and uniformly, then outputting their
field-pairing generator using the basis-expanded family, has every nonzero Fourier
coefficient of absolute value at most $ n/2^\ell$.
\end{theorem}

\begin{theorem}[Dyadic parameter identity]
\lean{FABL.basisExpandedSmallBias_dyadicParameter}
\label{FABL.basisExpandedSmallBias_dyadicParameter}
\leanok
\uses{FABL.dyadicFieldRatio_eq_invPow}
\difficulty{1}
Let $n$, $t$, and $\ell$ be natural numbers such that $t \le \ell$ and $n=2^{\ell-t}$.
Then, as real numbers, $n/2^\ell=(1/2)^t$.
\end{theorem}

\begin{theorem}[Bias of the basis-expanded generator]
\lean{FABL.basisExpandedSmallBiasGeneratorDensity_isBiased_dyadic}
\proofsource{Analysis-of-Boolean-Functions-by-Ryan-ODonnell}{
  pdf-ffbe9c7b3306-p153-b004,
  pdf-ffbe9c7b3306-p166-b005
}
\label{FABL.basisExpandedSmallBiasGeneratorDensity_isBiased_dyadic}
\leanok
\uses{FABL.ProbabilityDensity.IsBiased, FABL.basisExpandedSmallBiasGeneratorDensity, FABL.basisExpandedSmallBiasGeneratorDensity_isBiased, FABL.basisExpandedSmallBias_dyadicParameter}
\difficulty{3}
Let $\ell>0$ and let $t,n$ be nonnegative integers with $t\le\ell$ and $n=2^{\ell-t}$.
The basis-expanded small-bias generator density on $\mathbb{F}_2^{n\ell}$ has every
nonzero vector Fourier coefficient bounded in absolute value by $2^{-t}$.
\end{theorem}
